
\begin{proposition}[Credited re-derivation: Excluded algebraic dimensions]\label{cdp:st:E3}
A positive-dimensional compact Moishezon manifold has $b_2>0$.
A compact complex manifold homeomorphic to $S^6$ has algebraic
dimension $0$ or $1$; in the latter case its algebraic reduction is
holomorphic onto $\mathbb P^1$.
\end{proposition}

The Moishezon assertion is \citet[Lemma~1.4]{CDP98}. The sphere assertions follow from \citet[Theorem~3.1 and Corollary~3.2, p.~9]{LRS2019}. We retain the complete deduction as a known consequence.

\begin{proof}
For the first assertion take a projective modification
$\pi:\widehat Y\to Y$. Choose a general very ample divisor
$\widehat D$ and a general complete-intersection curve
$\widehat C$ in the smooth projective manifold $\widehat Y$ so that
they meet outside the exceptional locus and their intersection is
finite. Their images $D,C$ meet properly and nontrivially. Since $Y$
is smooth, $D$ is Cartier and its effective restriction to the
normalization of $C$ has positive degree. Thus
$\langle c_1(\mathcal O_Y(D)),[C]\rangle>0$, proving $b_2(Y)>0$.
This is the proof of CDP98, Lemma~1.4 in the opened arXiv source.

On a six-sphere the Euler number is $2$. Statement \ref{cdp:st:E2} consequently
excludes every meromorphic function with an indeterminate map to
$\mathbb P^1$. If two meromorphic functions were algebraically
independent, both would extend to holomorphic maps to
$\mathbb P^1$. Their product map to $\mathbb P^1\times\mathbb P^1$
has rank two at some point. In local target coordinates $u,v$
centred there, the meromorphic function $u/v$ pulled back to $Y$
has an indeterminacy: its restrictions to the local curves
$u=0$ and $v=0$ have distinct limits. This contradicts \ref{cdp:st:E2}.
Thus $a(Y)\leq1$. For $a(Y)=1$, the algebraic reduction cannot be
non-holomorphic by \ref{cdp:st:E2}, and its base has genus zero because pullback
injects $H^1(C,\mathbb Q)$ into $H^1(Y,\mathbb Q)=0$.
\end{proof}


